\documentclass[../main.tex]{subfiles}
\begin{document}
\subsection{質量傳輸定理推導}
\begin{itemize}
\item Governing Equation:
  \begin{itemize}
    \item 對一個物質的質量守恆:
    \begin{equation}
      m_{\text{in}} - m_{\text{out}} + m_{\text{gen}} = \Delta m_{\text{acc}} \label{eq:mass_conservation_general}
    \end{equation}
    故同之前熱傳與動量傳輸:
    \begin{equation}
      -\iint \vec{\bm n}_i \cdot \vec{\bm n}_sdS + \iiint r_i dV = \iiint \frac{\partial \rho_i}{\partial t} dV \label{eq:mass_conservation_integral}
    \end{equation}
    而其中$\vec{\bm n}_s$為控制體積表面的法向量\\
    同樣根據高斯定理，將表面積分轉成體積積分:
    \begin{align}
      & - \iiint \nabla \cdot \vec{\bm n}_i dV + \iiint r_i dV = \iiint \frac{\partial \rho_i}{\partial t} dV \nonumber\\ 
      \Rightarrow \quad & \iiint\left(
        \frac{\partial \rho_i}{\partial t} + \nabla \cdot \vec{\bm n}_i - r_i
      \right) dV = 0 \nonumber\\
      \Rightarrow \quad & \boxed{\frac{\partial \rho_i}{\partial t} + \nabla \cdot \vec{\bm n}_i  -r_i = 0} \\
      \text{or} \quad & \boxed{\frac{\partial \rho_i}{\partial t} + \nabla \cdot (\rho_i \vec{\bm u}) - r_i = 0} \label{eq:mass_conservation_differential}
    \end{align}
    \item 對所有物質的質量守恆:
    \begin{align}
      &\sum_i \frac{\partial \rho_i}{\partial t} + \sum_i\nabla \cdot (\rho_i \vec{\bm u}) - \cancelto{0}{\sum_i r_i} = 0 \nonumber\\
      &\boxed{\frac{\partial \rho}{\partial t}+ \nabla \cdot \rho \vec{\bm u}=0} \label{eq:total_mass_conservation}
    \end{align}
    而就會發現他就變成流力的Equation of Continuity了
    \item 寫成莫耳數的形式，有發生反應時會比較好計算
    \begin{equation}
      \frac{\partial C_i}{\partial t} + \nabla \cdot \vec{\bm N}_i - R_i = 0 \label{eq:molar_mass_conservation}
    \end{equation}
    $\vec{\bm N}_i=C_i\vec{\bm u}$\\
    P.S. $\sum_i R_i \neq 0$，會跟反應的化學計量有關
  \end{itemize}
  \item \fbox{Maxwell-Stefan Diffusion Equation}，Fick's Law前身，如何得到$J_i, j_i$:
  \begin{itemize}
    \item 想法:\\
    \fbox{物質$i$是受到所有非$i$的物質$j$的受力而移動}\\
    那些$j$施給$i$的力，會與$i$因為這個力施予其他$j$的力相等\\
    (牛頓第三運動定律，我被這些人擠=我擠這些人)
    \begin{equation}
      \vec{\bm \varphi}_i^j= k_{ij} \left(\vec{\bm u_j - \vec{\bm u_i} }\right)
    \end{equation}
    $\varphi_i^j$為$j$對$i$的單位體積下的作用力，$k_{ij}$為Friction Coefficient\\
    而$\vec{\bm u_j - \vec{\bm u_i} }$，則為兩物質的相對速度\\
    且$k_{ij}=k_{ji}$
    \item Friction Coefficient $k_{ij}$的物理意義:\\
    作用力的大小，與單位時間內的碰撞頻率有關\\
    而碰撞頻率，又跟兩物質的速度差成正比\\
    而Stefan-Maxwell假設$\boxed{k_{ij}\propto C_i C_j}$\\
    其比例常數取決於物質特性、溫度、以及整體濃度\\
    將物質特性部分寫作$\mathbcal{D}_{ij}$，有著\fbox{單位 $\frac{m^2}{s}$}
    \begin{equation}
      k_{ij} = \frac{C_iC_j RT}{C \mathbcal{D}_{ij}},\quad C=\sum_i C_i
    \end{equation}
    P.S. 同理$\mathbcal{D}_{ij}=\mathbcal{D}_{ji}$，稱之為\fbox{Maxwell-Stefan Diffusivity}
    \item 對物質$i$，受到所有$j$的作用力總和
    \begin{equation}
      \vec{\bm \varphi}_i = \sum_{j \neq i} \vec{\bm \varphi}_i^j = \sum_{j \neq i} k_{ij} \left(\vec{\bm u_j - \vec{\bm u_i} }\right)
    \end{equation}
    $\vec{\bm \varphi}_i$，是Total Drag Force per Unit Volume on Species $i$\\
    而這個由速度現象產生的作用力來源，應該來自濃度不同導致的\fbox{化學能}不同
    \begin{equation}
      \text{Driving Force/V} = \vec{\bm d}_i = -C_i \nabla \mu_i,\quad 
      \mu_i = \left(
        \frac{\partial G}{\partial n_i}
      \right)_{T,P,n_j}
    \end{equation}
    P.S.  以單位來想的話，因為功是力乘上距離，故對能量取梯度\\
    單位會是\fbox{Force / mol}，再乘上濃度的單位\fbox{mol / m$^3$}\\
    即可得到\fbox{Force / m$^3$}\\
    而因為Driving Force是物質感受周遭，而想要對外做功\\
    Drag Force則是周遭物質對物質$i$施加的阻力\\
    故兩者應該大小相等，方向相反
    \item Maxwell-Stefan Diffusion Equation:
    \begin{equation}
      \boxed{C_i \nabla \mu_i = \sum_{j\neq i}k_{ij}\left(
        \vec{\bm u_j - \vec{\bm u_i} }
      \right)} \label{eq:maxwell_stefan_diffusion}
    \end{equation}
    如果計算整體加總，會發現根據Gibbs-Duhem Equation
    \begin{equation}
      \boxed{
        SdT -VdP + \sum_i n_i d\mu_i = 0
      }
    \end{equation}
    在\fbox{等溫、等壓}下，後面那項$\sum_i n_i d\mu_i = 0$\\
    這也剛好是Maxwell-Stefan Diffusion Equation加總後左側的結果\\
    故
    \begin{equation}
      \sum_i C_i \nabla \mu_i = 0 \implies \sum_i \sum_{j\neq i}\frac{C_i C_j RT}{C \mathbcal{D}_{ij}} \left(
        \vec{\bm u_j - \vec{\bm u_i} }
      \right) = 0
    \end{equation}
    \item 對每個物質$i$，可以寫成以$\vec{\bm N}_i$的形式\\
    根據定義$\vec{\bm N}_i = C_i \vec{\bm u}_i$
    \begin{equation}
      \vec{\bm u}_i = \frac{\vec{\bm N}_i}{C_i}
    \end{equation}
    代入(\ref{eq:maxwell_stefan_diffusion})，可以消除$C_i, C_j$
    \begin{align}
      C_i \nabla \mu_i &= \sum_{j\neq i}k_{ij}\left(
        \vec{\bm u_j - \vec{\bm u_i} }
      \right)\nonumber\\
      &=\sum_{j\neq i} \frac{C_i C_j RT}{C \mathbcal{D}_{ij}} \left( 
        \vec{\bm u}_j - \vec{\bm u}_i
      \right) \nonumber\\
      &= \sum_{j\neq i} \frac{C_i C_jRT}{C \mathbcal{D}_{ij}} \left(
        \frac{\vec{\bm N}_j}{C_j} - \frac{\vec{\bm N}_i}{C_i}
      \right) \nonumber\\
      &= \sum_{j\neq i} \frac{C_i C_jRT}{C \mathbcal{D}_{ij}} \left(
        \frac{\vec{\bm N}_jC_i - \vec{\bm N}_iC_j}{C_iC_j}
      \right)\nonumber\\
      &= \sum_{j\neq i} \frac{ RT}{C\mathbcal{D}_{ij}}
      \left(\vec{\bm N}_jC_i - \vec{\bm N}_iC_j\right)
    \end{align}
    把整體總濃度$C$移入括號內$\frac{C_i}{C} = X_i$\\
    即為物質$i$的莫耳分率
    \begin{equation}
      \boxed{C_i \nabla \mu_i = \sum_{j\neq i} \frac{ RT}{\mathbcal{D}_{ij}}
      \left(x_i \vec{\bm N}_j - x_j \vec{\bm N}_i\right)} \label{eq:maxwell_stefan_diffusion_final}
    \end{equation}
    \item 理想氣體下的Maxwell-Stefan Diffusion Equation:
    \begin{equation}
      \mu_i = \mu_i^\circ + RT \ln (x_i P)
    \end{equation}
    P.S. 如果不是理想氣體，要將活度係數考慮進去，換掉$x_i$
    \begin{equation}
      \mu_i = \mu_i^\circ(T,P) + RT \ln (a_i) \label{eq:chemical_potential_non_ideal}
    \end{equation}
    而$\nabla \mu_i$為(定溫、定壓下)
    \begin{align}
      \nabla \mu_i &= \nabla \left(
        \mu_i^\circ + RT \ln( x_i P)
      \right) \nonumber\\
      &= RT \nabla \ln (x_i P) \nonumber\\
      &= RT \left(\frac{1}{x_i P}\cdot P \nabla x_i\right) \nonumber\\
      &= \frac{RT}{x_i} \nabla x_i \label{eq:chemical_potential_gradient_ideal_gas}
    \end{align}
    代回(\ref{eq:maxwell_stefan_diffusion_final})，可以消掉$RT$
    \begin{align}
      C_i \nabla \mu_i &= \sum_{j\neq i} \frac{ RT}{\mathbcal{D}_{ij}}
      \left(x_i \vec{\bm N}_j - x_j \vec{\bm N}_i\right) \nonumber\\
      \frac{C_i \cancel{RT}}{x_i} \nabla x_i &= \sum_{j\neq i} \frac{\cancel{RT}}{\mathbcal{D}_{ij}}
      \left(x_i \vec{\bm N}_j - x_j \vec{\bm N}_i\right) \nonumber\\
      \frac{C_i}{\frac{C_i}{C}} \nabla x_i &= \sum_{j\neq i} \frac{1}{\mathbcal{D}_{ij}}
      \left(x_i \vec{\bm N}_j - x_j \vec{\bm N}_i\right) \nonumber\\
      \implies \quad \nabla x_i &= \sum_{j\neq i} \frac{1}{C \mathbcal{D}_{ij}} \left(
        x_i \vec{\bm N}_j - x_j \vec{\bm N}_i
      \right) \label{eq:maxwell_stefan_diffusion_ideal_gas}
    \end{align}
    如果是二元系統，右邊的$(x_i \vec{\bm N}_j - x_j \vec{\bm N}_i)$可以合併成一個變數:
    \begin{align}
      x_A \vec{\bm N}_B - x_B \vec{\bm N}_A &= 
      \underbrace{x_A \vec{\bm N}_B}_A - \underbrace{x_A \vec{\bm N}_B}_B
      + \underbrace{x_A\vec{\bm N}_A}_A -\underbrace{ x_A \vec{\bm N}_A}_B \nonumber\\
      &= x_A \left(
        \vec{\bm N}_A + \vec{\bm N}_B
      \right) - \vec{\bm N}_A \left(
        x_A + x_B
      \right) \nonumber\\
      &= x_A \left(\vec{\bm N}_A + \vec{\bm N}_B\right) - \vec{\bm N}_A
    \end{align}
    \item 二元系統、理想氣體，定溫定壓下的Stefan-Maxwell Diffusion Equation\\
    物質為A,B，以A來看
    \begin{align}
      C_A \nabla \mu_A &= \frac{RT}{\mathbcal{D}_{AB}} (x_A \vec{\bm N}_B - x_B \vec{\bm N}_A) \nonumber\\
      & = \frac{RT}{\mathbcal{D}_{AB}} \left(x_A \left(\vec{\bm N}_A + \vec{\bm N}_B\right) - \vec{\bm N}_A\right) \\
      \nabla x_A &= \frac{1}{C \mathbcal{D}_{AB}} \left(
        x_A \left(\vec{\bm N}_A + \vec{\bm N}_B\right) - \vec{\bm N}_A
      \right)
    \end{align}
    \item 二元系統，非理想流體，非定溫定壓的Stefan-Maxwell Diffusion Equation\\
    物質為A,B，以A來看\\
    化學能的關係式為(\ref{eq:chemical_potential_non_ideal})
    \begin{equation}
      \mu_i = \mu_i^\circ(T,P) + RT \ln (a_i) 
    \end{equation}
    而對$\mu_A(T,p,a_A)$做全微分，要對溫度、壓力、活度做偏微分
    \begin{align}
      d\mu_A &= \left(\frac{\partial \mu_A}{\partial T}\right)_{P,a_A} dT +
      \left(\frac{\partial \mu_A}{\partial P}\right)_{T,a_A} dP +
      \left(\frac{\partial \mu_A}{\partial a_A}\right)_{T,P} da_A \nonumber\\
      &= -S_A dT + V_A dP + RT \frac{1}{a_A} da_A
    \end{align}
    $S$為物質A的Molar Entropy，$V$為Molar Volume\\
    而對空間取梯度
    \begin{equation}
      \nabla \mu_A = -S_A \nabla T + V_A \nabla P + RT \nabla\ln (a_A)
    \end{equation}
    而前兩項可以用狀態方程算出來，移到左側就修正了\\
    故假設之後提到的$\nabla \mu_i$，都是減去這兩項之後的結果\\
    而因為等號右邊是用$x_i$表示，因此要將$a_i$換成$x_i$
    \begin{equation}
      \nabla \mu_A +S_A \nabla T - V_A \nabla P = RT \left(
        \frac{\partial (\ln a_A)}{\partial x_A}
      \right)_{T,P} \nabla x_A \label{eq:chemical_potential_gradient_non_ideal}
    \end{equation}
    代回(\ref{eq:maxwell_stefan_diffusion_final})
    \begin{equation}
      C_A RT \left(
        \frac{\partial (\ln a_A)}{\partial x_A}
      \right)_{T,P} \nabla x_A = \frac{RT}{\mathbcal{D}_{AB}} \left(
        x_A \left(\vec{\bm N}_A + \vec{\bm N}_B\right) - \vec{\bm N}_A
      \right)
    \end{equation}
    可以消掉$RT$
    \begin{equation}
      C_A \left(
        \frac{\partial (\ln a_A)}{\partial x_A}
      \right)_{T,P} \nabla x_A = \frac{1}{\mathbcal{D}_{AB}} \left(
        x_A \left(\vec{\bm N}_A + \vec{\bm N}_B\right) - \vec{\bm N}_A
      \right)
    \end{equation}
    \item 移項，用\fbox{$\vec{\bm N}_A$}來表示，會得到\fbox{Fick's Law}
    \begin{align}
      C_A\mathbcal{D}_{AB}\left(
        \frac{\partial (\ln a_A)}{\partial x_A}
      \right)_{T,P} \nabla x_A &= x_A \left(\vec{\bm N}_A + \vec{\bm N}_B\right) - \vec{\bm N}_A \nonumber\\
      \vec{\bm N}_A &= x_A \left(\vec{\bm N}_A + \vec{\bm N}_B\right) - \boxed{C_A}\mathbcal{D}_{AB}\left(
        \frac{\partial (\ln a_A)}{\partial x_A}
      \right)_{T,P} \nabla x_A \nonumber\\
      \vec{\bm N}_A &= x_A \left(\vec{\bm N}_A + \vec{\bm N}_B\right) - \boxed{x_AC}\mathbcal{D}_{AB}\left(
        \frac{\partial (\ln a_A)}{\partial x_A}
      \right)_{T,P} \nabla x_A \nonumber\\
      \vec{\bm N}_A &= x_A \left(\vec{\bm N}_A + \vec{\bm N}_B\right) -C \mathbcal{D}_{AB}\left(
        \frac{\partial (\ln a_A)}{\partial \ln(x_A)}
      \right)_{T,P} \nabla x_A
    \end{align}
    定義\fbox{Fickian Diffusivity}為
    \begin{equation}
      \boxed{D_{AB} = \mathbcal{D}_{AB} \left(
        \frac{\partial (\ln a_A)}{\partial \ln(x_A)}
      \right)_{T,P}}
    \end{equation}
    \fbox{$\mathbcal{D}_{AB}$ 比起 $D_{AB}$ 跟濃度更無關}\\
    最後得到Fick's 1st Law的形式
    \begin{equation}
      \boxed{\vec{\bm N}_A = x_A \left(\vec{\bm N}_A + \vec{\bm N}_B\right) - C D_{AB} \nabla x_A} \label{eq:ficks_law_general}
    \end{equation}
  \end{itemize}
  \item 從Fick's 1st Law (\ref{eq:ficks_law_general})推出的其他形式:
  \begin{itemize}
    \item 用莫耳平均速度$\vec{\bm u}^\ast$來表示\\
    (大寫可加成，$J_i$會消掉)
    \begin{equation}
      \sum N_i = \vec{\bm u}^\ast \sum C_i
    \end{equation}
    將$x_A \left(\vec{\bm N}_A + \vec{\bm N}_B\right)$代入莫耳平均速度的定義(\ref{eq:molar_flux_decomposition_volume_average})\
    \begin{equation}
      x_A \left(\vec{\bm N}_A + \vec{\bm N}_B\right) = x_A\left(
        C_A \vec{\bm u}_A + C_B \vec{\bm u}_B
      \right) = x_A C \vec{\bm u}^\ast = C_A \vec{\bm u}^\ast
    \end{equation}
    代回(\ref{eq:ficks_law_general})
    \begin{equation}
     \boxed{ \vec{\bm N}_A = C_A \vec{\bm u}^\ast - \underbrace{C D_{AB} \nabla x_A}_{\vec{\bm J}_A}}
    \end{equation}
    回到最一開始的Governing Equation(\ref{eq:molar_flux_decomposition_volume_average})\\
    \fbox{第一項就是對流項，第二項是分子擴散項}
    \item 如果用質量分率來表示Fick's Law:
    \begin{equation}
      \vec{\bm  n}_A = \rho_A \vec{\bm u}^\ast + \vec{\bm j}_A
    \end{equation}
    我們知道$C \vec{\bm u}^\ast = \sum_i C_i \vec{\bm u}_i$\\
    乘上$M_A$，並利用(\ref{eq:species_molar_concentration})，可以得到
    \begin{equation}
      \rho \vec{\bm u}^\ast = \sum_i \rho_i \vec{\bm u}_i
    \end{equation}
    兩者個連接差在一個是以莫耳數表示，一個是以質量表示
    \begin{equation}
      \vec{\bm  N}_A  = \frac{\vec{\bm n}_A}{M_A} = \frac{1}{M_A}\left(
        \rho_A \vec{\bm u} + \vec{\bm j}_A
      \right) =  C_A \vec{\bm u}^\ast - \vec{\bm J}_A^\ast
    \end{equation}
    而$j_A=\rho_A(\vec{\bm u}_A -\vec{\bm_u})$，是相對於體積平均速度的分子擴散\\
    將體積平均速度展開:
    \begin{equation}
      \vec{\bm u} = \frac{\left(\rho_A \vec{\bm u}_A + \rho_B \vec{\bm u}_B\right)}{\rho_A + \rho_B}
    \end{equation}
    定義
    \begin{equation}
      \omega_i = \frac{\rho_i}{\sum_j \rho_j}
    \end{equation}
    為物質$i$的質量分率\\
    則
    \begin{equation}
      \vec{\bm u} = \omega_A \vec{\bm u}_A + \omega_B \vec{\bm u}_B
    \end{equation}
    代入$j_A$中
    \begin{align}
      \vec{\bm j}_A &= \rho_A \left(\vec{\bm u}_A - \vec{\bm u}\right) \nonumber\\
      &= \rho_A \left(\vec{\bm u}_A - \omega_A \vec{\bm u}_A - \omega_B \vec{\bm u}_B\right) \nonumber\\
      &= \rho_A \left(
        (1-\omega_A) \vec{\bm u}_A - \omega_B \vec{\bm u}_B
      \right) \nonumber\\
      &= \rho_A \left(
        \omega_B \vec{\bm u}_A - \omega_B \vec{\bm u}_B
      \right) \nonumber\\
      &= \rho_A \omega_B \left(
        \vec{\bm u}_A - \vec{\bm u}_B
      \right)
    \end{align}
    而相似的，大寫的分子擴散
    \begin{align}
      \vec{\bm J}_A &= C_A \left(\vec{\bm u}_A - \vec{\bm u}^\ast\right) \nonumber\\
      &= C_A \left(\vec{\bm u}_A - x_A \vec{\bm u}_A - x_B \vec{\bm u}_B\right) \nonumber\\
      &= C_A \left(
        (1-x_A) \vec{\bm u}_A - x_B \vec{\bm u}_B
      \right) \nonumber\\
      &= C_A \left(
        x_B \vec{\bm u}_A - x_B \vec{\bm u}_B
      \right) \nonumber\\
      &= C_A x_B \left(
        \vec{\bm u}_A - \vec{\bm u}_B
      \right)
    \end{align}
    因此兩者之間具有以下比例關係
    \begin{equation}
      \boxed{\frac{\vec{\bm j}_A}{\vec{\bm J}_A} = \frac{\rho_A \omega_B}{C_A x_B} = \frac{M_A \omega_B}{x_B}
      } \label{eq:molar_to_mass_diffusion_flux}
    \end{equation}
    另外$x$和$\omega$之間的關係為
    \begin{equation}
      x_A = \frac{\frac{\omega_A}{M_A}}{\frac{\omega_A}{M_A} + \frac{\omega_B}{M_B}}
    \end{equation}
    而因為$\omega_B = 1 - \omega_A$，因此
    \begin{equation}
      x_A = \frac{\frac{\omega_A}{M_A}}{\frac{\omega_A}{M_A} + \frac{1-\omega_A}{M_B}}
    \end{equation}
    微分:
    \begin{equation}
      \frac{d}{d\omega_A} x_A =\frac{
        \frac{1}{M_A} \left(
          \frac{\omega_A}{M_A} + \frac{1-\omega_A}{M_B}
        \right) - \frac{\omega_A}{M_A} \left(
          \frac{1}{M_A} - \frac{1}{M_B}
        \right)
      }{\left(
        \frac{\omega_A}{M_A} + \frac{1-\omega_A}{M_B}
      \right)^2}
    \end{equation}
  令平均分子量的倒數
  \begin{equation}
    \frac{1}{\overline{M}} = \frac{\omega_A}{M_A} + \frac{1-\omega_A}{M_B}
  \end{equation}
  則
  \begin{equation}
    \frac{d}{d\omega_A} x_A = \frac{1}{\overline{M}^2} \cdot \frac{1}{M_A M_B}
  \end{equation}
  因此
  \begin{equation}
    \frac{d x_A}{d \omega_A} = \frac{\overline{M}^2}{M_A M_B} \label{eq:dxA_domegaA}
  \end{equation}
  故
  \begin{equation}
    \nabla x_A = \frac{\overline{M}^2}{M_A M_B} \nabla \omega_A \label{eq:grad_xA_to_grad_omegaA}
  \end{equation}
  於是Fick's Law可以寫成以質量分率表示
  \begin{align}
    \vec{\bm J}_A &= - C D_{AB} \nabla x_A \nonumber\\
    &= - C D_{AB} \cdot \frac{\overline{M}^2}{M_A M_B} \nabla \omega_A 
  \end{align}
  由(\ref{eq:molar_to_mass_diffusion_flux})，可得
  \begin{equation}
    \vec{\bm j}_A = \frac{M_A \omega_B}{x_B} \vec{\bm J}_A
  \end{equation}
  代入:
  \begin{align}
    \vec{\bm j}_A &= -CD_{AB}\cdot\left(
      \frac{M_A \omega_B}{x_B}
    \right)\cdot \frac{\overline{M}^2}{M_A M_B} \nabla \omega_A \nonumber\\
    &= -CD_{AB} \cdot \frac{\omega_B}{x_B} \cdot \frac{\overline{M}^2}{M_B} \nabla \omega_A
  \end{align}
  而$C\cdot \overline{M} = \rho$，$x_B = \omega_B\frac{\overline{M}}{M_B}$\\
  故
  \begin{align}
    \vec{\bm j}_A &= - CD_{AB} \cdot \frac{\omega_B}{x_B} \cdot \frac{\overline{M}^2}{M_B} \nabla \omega_A \nonumber\\
    &= - \underbrace{C\overline{M}}_{\rho} D_{AB}\cdot \left(
      \underbrace{\omega_B\frac{\overline{M}}{M_B}}_{x_B}\cdot \frac{1}{x_B}
    \right) \nabla \omega_A \nonumber\\
    &= - \rho D_{AB} \nabla \omega_A
  \end{align}
  因此以質量分率表示的Fick's Law為
  \begin{equation}
    \boxed{\vec{\bm j}_A = - \rho D_{AB} \nabla \omega_A} \label{eq:ficks_law_mass_fraction}
  \end{equation}
  或者:
  \begin{equation}
    \vec{\bm n}_A = \rho_A \vec{\bm u}^\ast - \rho D_{AB} \nabla \omega_A
  \end{equation}
  注意這只有二元系統可以這樣換
  \item 簡略形式，二元系統，密度固定，Fickian擴散係數固定，能量平衡式
  \begin{align}
    \frac{\partial \omega_A}{\partial t} + \nabla \cdot \vec{\pm n}_A -r_A &= 0 \nonumber\\
    \frac{\partial \omega_A}{\partial t} + \nabla \cdot \left(
      \rho_A \vec{\bm u} + \vec{\bm j}_A
    \right) - r_A &= 0 \nonumber\\
    \frac{\partial \omega_A}{\partial t} + \nabla \cdot \left(
      \rho_A \vec{\bm u} - \rho D_{AB} \nabla \omega_A
    \right) - r_A &= 0 \nonumber\\
    \frac{\partial \omega_A}{\partial t} + \nabla \cdot \left(
      \rho_A \vec{\bm u} -  D_{AB} \nabla \rho_A
    \right) - r_A &= 0 \nonumber\\
    \frac{\partial \omega_A}{\partial t} + \rho_A \cancelto{0}{\nabla \cdot \vec{\bm u}} 
    + \vec{\bm u} \cdot \nabla \rho_A -D_{AB} \nabla^2 \rho_A - r_A &=0 \nonumber\\
    (\div M_A)\implies \quad \boxed{\frac{\partial C_A}{\partial t} 
    + \vec{\bm u} \cdot \nabla C_A - D_{AB} \nabla^2 C_A - R_A} &=0\label{eq:mass_transport_simplified }
  \end{align}
  \item 更簡略形式，二元系統，密度固定，擴散係數固定，一物濃度極低\\
  令$C_i\ll C_s$，$C_i$是溶質，$C_s$是溶劑\\
  因此這時$C=C_s + C_i \approx C_s$為常數\\
  平均速度也只會是溶劑的，$\vec{\bm u} \approx \vec{\bm u}_s$\\
  化學能變化:
  \begin{align}
    \mu_i &= \mu_i^\circ + RT \ln (x_i P) \nonumber\\
    \nabla \mu_i &= RT \nabla (\ln x_i )
  \end{align}
  從(\ref{eq:maxwell_stefan_diffusion})開始
  \begin{align}
    C_i \nabla \mu_i &= RT\frac{C_i C_s}{C \mathbcal{D}_{is}} \left(
       \vec{\bm u}_s - \vec{\bm u}_i
    \right) \nonumber\\
    &\approx RT \frac{C_i C_s}{C_s \mathbcal{D}_{is}} \left(
       \vec{\bm u} - \vec{\bm u}_i
    \right) \nonumber\\
    C_i RT\nabla (\ln C_i) &= RT \frac{C_i}{\mathbcal{D}_{is}} \left(
       \vec{\bm u} - \vec{\bm u}_i
    \right) \nonumber\\
    \nabla (\ln C_i) &= \frac{1}{\mathbcal{D}_{is}} \left(
       \vec{\bm u} - \vec{\bm u}_i
    \right) \nonumber\\
    \frac{1}{C_i}\nabla C_i &= \frac{1}{\mathbcal{D}_{is}} \left(
       \vec{\bm u} - \vec{\bm u}_i
    \right) \nonumber\\
    -\mathbcal{D}_{is} \nabla C_i &= C_i\vec{\bm u} - \underbrace{C_i \vec{\bm u}_i}_{\vec{\bm N}_i} \nonumber\\
    \vec{\bm N}_i &= C_i \vec{\bm u} - \underbrace{\mathbcal{D}_{is} \nabla C_i}_{\vec{\bm J}_i} \nonumber\\
    \implies \quad &\boxed{\vec{\bm J}_i = -\mathbcal{D}_{is} \nabla C_i} \label{eq:ficks_law_dilute_solution}
  \end{align}
  能量平衡式
  \begin{align}
    \frac{\partial C_i}{\partial t} + \nabla \cdot \vec{\bm N}_i - R_i &= 0 \nonumber\\
    \frac{\partial C_i}{\partial t} + \vec{\bm u} \cdot \nabla C_i - \mathbcal{D}_{is} \nabla^2 C_i - R_i &= 0
  \end{align}
  假設穩態，沒有反應
  \begin{equation}
    \boxed{\vec{\bm u} \cdot \nabla C_i = \mathbcal{D}_{is} \nabla^2 C_i} \label{eq:mass_transport_dilute_solution_steady_state}
  \end{equation}
    \item 對於incomprssible flow，$\nabla\cdot \vec {\bm u} = 0$，則
  \begin{equation}
    \frac{\partial C_A}{\partial t} + \vec {\bm u} \cdot \nabla C_A = D_{AB}\nabla^2 C_A +R_A
  \end{equation}
  或寫成:
  \begin{equation}
    \frac{DC_A}{Dt} = D_{AB}\nabla^2 C_A +R_A
  \end{equation}
  \item 對於incomprssible flow，且沒有質量生成$R_A=0$，則
  \begin{equation}
    \frac{\partial C_A}{\partial t} + \vec {\bm u} \cdot \nabla C_A = D_{AB}\nabla^2 C_A
  \end{equation}
  或寫成:
  \begin{equation}
    \frac{DC_A}{Dt} = D_{AB}\nabla^2 C_A
  \end{equation}
  \item \fbox{Transient Diffusion}對於無流體流動$\vec {\bm u}=0$，且沒有反應
  \begin{equation}
    \frac{\partial C_A}{\partial t} = D_{AB}\nabla^2 C_A
  \end{equation}
  此為Fick's Second Law
  \item Transient Diffusion in a slab，表面濃度固定，擴散係數固定
  \begin{equation}
    \frac{C_s-\overline{C}}{C_s - C_0} = \frac{8}{\pi^2}\left[
      e^{-a_1\text{Fo}_m} + \frac{1}{9}e^{-9a_1\text{Fo}_m} + \frac{1}{25}e^{-25a_1\text{Fo}_m} + \cdots
    \right]
  \end{equation}
  其中$s$是one-half slab的厚度，$a_1=(\pi/2)^2$是一個幾何/特徵相關的常數\\
  $C_0$是初始濃度，$\overline{C}$是Slab內的平均濃度，$C_s$是Slab表面的濃度\\
  $\text{Fo}_m = \frac{D_{AB} t}{s^2}$為質傳的Fourier Number
  \item 若為穩態，且無流體流動$\vec {\bm u}=0$，且沒有反應:
  \begin{equation}
    \nabla^2 C_A = 0
  \end{equation}
  此為Laplace Equation
\end{itemize}
\item Fick's Law的推導
\begin{figure}[H]
  \centering
  \begin{tikzpicture}[>=Latex, line cap=round, line join=round, thick]
    \draw (0,0) -- (0,3);
    \draw[<->] (-1.5, 1.5) -- (0, 1.5) node[midway, above] {$l_d$};
    \draw[<->] (0, 1.5) -- (1.5, 1.5) node[midway, above] {$l_d$};
    \draw[dashed] (-1.5, 0) -- (-1.5, 3);
    \draw[dashed] (1.5, 0) -- (1.5, 3);
    \node[anchor=south] at (-0.75, 3) {$n'$};
    \node[anchor=south] at (0.75, 3) {$n$};
    \draw[dashed, ->] (-2.2, 0.5) -- (-0.7, 0.5) node[right] {x};
    \draw[dashed, ->] (-0.5, 1) -- (1,1) node[right] {o};
  \end{tikzpicture}
  \caption{Fick's Law的推導示意圖}
\end{figure}
假設有一薄膜，左邊單位面積有$n'$個分子，右邊單位面積有$n$個分子\\
形成一個濃度梯度$\frac{dn}{dx}$，造成每秒鐘有$\kappa$個分子跳來跳去\\
在$x$方向上(擴散方向)，平均跳動$l_d$的距離\\
而能夠穿越的就只有膜兩側$l_d$距離內的分子
\begin{align}
  J_x &= \kappa n - \kappa n' = \kappa (n- n')  \\
  n' &= n + \frac{dn}{dx} l_d ,\quad \text{能跳過去的}\\
  \implies J_x &= -\kappa l_d \frac{dn}{dx} 
\end{align}
將單位面積的分子數$n$換成濃度$C$，則
\begin{equation}
  C = \frac{n}{l_d}
\end{equation}
因此
\begin{equation}
  J_x = -\kappa l_d^2 \frac{dC}{dx} \label{eq:ficks_law_derivation}
\end{equation}
若定義\fbox{Fickian Diffusivity}為
\begin{equation}
  D = \kappa l_d^2
\end{equation}
則就得到了Fick's Law的形式
\begin{equation}
  \boxed{J_x = -D \frac{dC}{dx}}
\end{equation}
\item 經驗式\\
  對比熱量輸送的經驗式$h$，因為氣體或液體的質傳行為複雜\\
  因此在分子擴散外的，質傳對流有以下Governing Equation，及經驗式
  \begin{itemize}
    \item 微觀下多分子的能量平衡式:\\
    Multicomponent Microscopic Energy Balance\\
    會連結質量傳輸的部分
    \begin{equation}
      \underbrace{\frac{\partial}{\partial t}\left(
        \rho \hat U + \frac{1}{2}\rho\left|\vec{\bm u}\right|^2
      \right)}_{E_{\text{acc}}} = \underbrace{-\nabla\cdot \vec{\bm e}}_{E_{\text{in}}-E_{\text{out}}}
      + \underbrace{\rho\left(\vec{\bm v}\cdot \vec{\bm g}\right)}_{E_{\text{gen}}}
    \end{equation}
    其中$\vec{\bm e}$為總能量通量:
    \begin{equation}
      \boxed{\vec{\bm e} = \left(\rho \hat U + \frac{1}{2}\rho\left|\vec{\bm u}\right|^2\right)\vec{\bm u} 
      + \vec{\bm q} + \underbrace{ \overline{\overline{\bm \pi}} \cdot \vec{\bm u}}_{P\bm \delta\bm\delta 
      + \overline{\overline{\bm \tau}}\cdot \vec{\bm u}}}
    \end{equation}
    對單物質系統
    \begin{equation}
      \vec{\bm q} = -k\nabla T
    \end{equation}
    對多物質系統
    \begin{equation}
      \vec{\bm q} = -k\nabla T + \sum_{i=1}^N \hat H_i \vec{\bm J}_i + \vec{\bm q}^{(T)}
    \end{equation}
    其中
    \begin{align}
      \hat H_i &= \left(\frac{\partial H}{\partial n_i}\right)_{T,P,n_j,j\neq i}, \quad \text{Molar Enthalpy of species }i\\
      \vec{\bm J}_i &= \text{Diffusive Molar Flux of species }i\\
      \vec{\bm q}^{(T)} &= \text{Duforer, Thermal Diffusion soret effect}
    \end{align}
    所以$\vec{\bm e}$又可以寫成其他形式
    \begin{equation}
      \boxed{\vec{\bm e} = \underbrace{\rho \hat U\vec{\bm u}}_{
        \sum (C_i\hat H_i) \vec{\bm u}
      } + \frac{1}{2}\rho\left|\vec{\bm u}\right|^2\vec{\bm u}
      + \left(-k\nabla T + \sum_{i=1}^N \hat H_i \vec{\bm J}_i + \vec{\bm q}^{(T)}\right)
      + P\bm \delta\bm\delta \cdot \vec{\bm u} + \overline{\overline{\bm \tau}}\cdot \vec{\bm u}}
    \end{equation}
    假設省略動能項$\frac{1}{2}\rho\left|\vec{\bm u}\right|^2\vec{\bm u}$
    、Duforer effect$\vec{\bm q}^{(T)}$\\
    耗散項$\overline{\overline{\bm \tau}} : \left(\nabla \vec{\bm u}\right)$、
    變形項$P\bm \delta\bm\delta \cdot \vec{\bm u}$\\
    則微觀能量方程式變成:
    \begin{align}
      \vec {\bm e} &= -k\nabla T + \sum_{i=1}^N \hat H_i \vec{\bm J}_i + \sum_i (C_i \hat H_i) \vec{\bm u}\\
      \quad  &= -k\nabla T +\sum_{i=1}^N\left[
        \hat H_i \left(\vec{\bm J}_i + C_i \vec{\bm u}\right)
      \right] \nonumber\\
      &= \boxed{-k\nabla T + \sum_{i=1}^N \hat H_i \vec{\bm N}_i}
    \end{align}
    若沒有重力，穩態下的微觀能量平衡，就跟Equation of Continuity一樣
    \begin{equation}
      \boxed{\nabla \cdot \vec{\bm e} = 0}
    \end{equation}
    \item  對流質量傳輸產生的條件:
      \begin{enumerate}
        \item 某成分有濃度差
        \item 有流動的連續介質
      \end{enumerate}
    \item 二元對流的質量傳輸
      \begin{equation}
        \boxed{\vec{\bm N}_A = k_c\cdot \Delta C_A = k_c\left(C_{A_1}-C_{A_2}\right)=k_c\cdot C\cdot\left(y_{A_1}-y_{A_2}\right)}
      \end{equation}
      $k_c$是質量傳輸係數\\
      也可以寫成壓力單位下:
      \begin{equation}
        \boxed{\vec{\bm N}_A = k_c\Delta C_A = k_p\Delta p_A = k_y\Delta y_A}
      \end{equation}
      $k_p$是壓力單位下的傳輸係數，$k_y$是莫爾單位下的傳輸係數\\
      P.S. $k_y = k_p\cdot P = k_c C$\\
      同熱對流的$h$，質量傳輸係數$k_{c,p,y}$也很難獲得\\
      因為會受容器的形狀、流速、流動方式、流體的物性(密度、黏度、熱容、熱導係數)、質傳方式而影響\\
      通常使用質量傳輸經驗式，計算，再代入獲得$k_{c,p,y}$
    \item 而對比熱對流會用Nusselt Number或Stanton Number (Heat) 來求出$h$\\
      質傳會用Sherwood Number或Stanton Number (Mass)來求出$k_c,k_y,k_p$
  \end{itemize}
\item 小結，解題分類:
\begin{itemize}
  \item 多為二元系統，一維、穩態\\
  不是就乖乖搞Maxwell-Stefan Diffusion(\ref{eq:maxwell_stefan_diffusion})
  \begin{equation}
    C_A \nabla \mu_A = \frac{RT}{\mathbcal{D}_{AB}} \left(
      x_A \left(\vec{\bm N}_A + \vec{\bm N}_B\right) - \vec{\bm N}_A
    \right)
  \end{equation}
  \item 若為固體，則可以使用Fick's Law
  \begin{equation}
    \vec{\bm J}_A = - D_{AB} \nabla C_A
  \end{equation}
  \item 若為流體，則分子擴散+質傳對流，可分為以下幾種情形\\
  可以說是簡單的質傳問題
  \begin{enumerate}
    \item 沒有反應，且一個物質靜止(Stagnant)
    \begin{equation}
      \vec{\bm N}_B = 0 \implies \vec{\bm N}_A = - C D_{AB} \nabla x_A
    \end{equation}
    \item 沒有反應，A和B是等莫耳交換流(由一方驅動，另一方填補)
    \begin{equation}
      \vec{\bm N}_A = -\vec{\bm N}_B \implies \vec{\bm N}_A = \frac{-C D_{AB}}{1-2x_A} \nabla x_A
    \end{equation}
    \item 沒有反應，A和B其中一方非常稀薄(Dilute Solution)
    \begin{equation}
      x_A \to 0 \quad \text{or}\quad x_B \to 0 \implies \vec{\bm N}_A = - C D_{AB} \nabla x_A
    \end{equation}
    \item 沒有反應\\
    使用Fick's First Law的完整形式
    \begin{equation}
      \vec{\bm N}_A = x_A \left(\vec{\bm N}_A + \vec{\bm N}_B\right) - C D_{AB} \nabla x_A
    \end{equation}
    \item 有反應，但反應是非勻相反應(Heterogeneous Reaction)
    \item 有反應，且反應是勻相反應(Homogeneous Reaction)
    \item 有反應，但是是孔洞觸媒的反應\\
    以上這幾種的解題，又分成短時間與長時間的質傳問題\\
    兩者都算是穩態，但長時間則變成假的穩態(Quasi-Steady State)\\
    和必須要用非穩態算的不同，簡單的質傳問題\\
    會假設擴散需要的時間遠小於巨觀變化(例如液面下降)所需時間\\
    長時間的問題，只是在算完擴散項後，再帶入巨觀變化的方程式中\\
    而不是直接把兩個變化同時算在一起
    \item 雖然是二維，但沒有反應、穩態
    \begin{equation}
      \frac{\partial^2 C_A}{\partial x^2} + \frac{\partial^2 C_A}{\partial y^2} = 0
    \end{equation}
    \item 一維、沒有反應，但不是穩態
    \begin{equation}
      \frac{\partial C_A}{\partial t} = D_{AB} \frac{\partial^2 C_A}{\partial x^2}
    \end{equation}
  \end{enumerate}
\end{itemize}
\end{itemize}
\end{document}